\section{Recompensas automáticas y señales de comprensión}
\subsection{Representación del objetivo con imágenes y lenguaje}
El requisito central de las recompensas autónomas es reducir el trabajo de ingeniería manual. Chelsea describió el pipeline para automatizar la recompensa con una \textbf{imagen objetivo} o una \textbf{descripción en lenguaje}: un modelo preentrenado extrae la semántica visual, se compara el escalamiento de los embedding del estado actual y del objetivo, y así se obtiene una recompensa de tipo arrastre.

\begin{figure}[H]
\centering
\includegraphics[width=0.9\textwidth]{slides-images/slide-02.jpg}
\caption{La Slide 2 muestra cómo la distancia entre los embedding del objetivo y del estado actual funciona como reward inmediata.}
\end{figure}

\begin{importantbox}{Embedding Distance Reward}
\begin{itemize}
\item Se usan modelos como VAE/CLIP para extraer `f(I)`;\\
\item La recompensa es $r_t = -\|\phi(I_t) - \phi(I_g)\|_2$;\\
\item Las diferencias de dominio pueden atenuarse con scratch/finetune, pero lo esencial es comparar estados mediante la representation.
\end{itemize}
\end{importantbox}

\subsection{Combinación de reward visual y de lenguaje}
La Slide 03 muestra el flujo que combina descripciones en lenguaje natural con embedding visuales para generar la reward: el LLM genera el task prompt, el VLM evalúa el current frame con base en el prompt y, al final, se filtra con un ensemble output. Con unos cuantos cases especiales de human-in-the-loop, el sistema puede distinguir las reward signal inseguras.

\begin{figure}[H]
\centering
\includegraphics[width=0.85\textwidth]{slides-images/slide-03.jpg}
\caption{La Slide 3 muestra el sistema en el que un LLM y un VLM se combinan para producir la reward, y subraya la importancia de la etapa de filter.}
\end{figure}

\begin{knowledgebox}{LLM + VLM Reward Pipeline}
\begin{itemize}
\item El LLM recibe el task prompt y genera una structured description;\\
\item El VLM (por ejemplo, CLIP/BLIP) califica el frame actual;\\
\item Un confidence filter (por ejemplo, la softmax temperature) controla la intensidad de la reward.
\end{itemize}
\end{knowledgebox}

\subsection{Análisis detallado del Reward Pipeline}
A partir del flujo de la slide 03, el expositor detalló las tres etapas del reward pipeline: \textbf{signal extraction} (obtener imágenes/lenguaje de los sensors), \textbf{affinity scoring} (embedding similarity) y \textbf{confidence gating} (filtrado por confianza y ensemble). Estas tres etapas deben modificarse de forma sincronizada después de cada actualización del modelo para garantizar que la reward not drift.

\begin{importantbox}{Marco abstracto del pipeline}
\begin{itemize}
\item Signal extraction: logging sensors + normalization;\\
\item Affinity scoring: distancia entre embedding o classifier probability;\\
\item Confidence gating: usar un filtro de threshold/ensemble/temporal smoothing para descartar las reward anómalas.
\end{itemize}
\end{importantbox}

\subsection{Autosupervisión basada en clasificadores de éxito}
El expositor recurre a anotaciones few-shot para que el robot aprenda a reconocer estados de éxito: se reúnen unas cuantas instantáneas positive/negative, se entrena un clasificador binario $\sigma(s)$ y la salida de ese classifier es la función de reward. Este enfoque es muy adecuado para tareas en las que «se desconoce el número de objetivos».

\begin{knowledgebox}{Success Classifier Pipeline}
\begin{itemize}
\item Reunir de 50 a 100 frames de «éxito» definidos por el usuario y generar muestras negativas;\\
\item Entrenar una CNN/RNN ligera que produzca probabilidades;\\
\item Usar la probabilidad directamente como reward, o combinarla con un decaimiento temporal para evitar una convergencia prematura.
\end{itemize}
\end{knowledgebox}

\subsection{Filtrado de reward multimodal}
Chelsea señaló además que, cuando CLIP/LLM generan la reward, es indispensable añadir un \textbf{filtrado por confiabilidad} para evitar que el modelo asigne puntajes altos a irrelevant features (como el fondo). Las prácticas habituales son hacer un ensemble de varias modalidades y agregar una regularización de temporal consistency.

\begin{warningbox}{No dejes que los resultados de hallucination del modelo pretrained lleguen a la policy}
Una reward basada en CLIP puede reforzar destellos del fondo. Se recomienda:
\begin{itemize}
\item Usar temporal smoothing para asegurar que la reward sea continua;\\
\item En los juicios multimodales, emplear majority vote/variance penalty;\\
\item Conservar un canal de audit human-in-the-loop para hacer override en caso de emergencia.
\end{itemize}
\end{warningbox}

\subsection{Resumen de la sección}
Las recompensas automáticas sustituyen la reward manual tradicional con diversas señales, como imágenes, lenguaje y clasificadores basados en modelos, y subrayan la necesidad de controlar la confianza de las salidas multimodales.
\newpage
